
\subsubsection{5.1 Main AUROC Results}

Table 1 reports AUROC values and 95\% bootstrap confidence intervals for all functional UQ methods on Llama-3-8B-Base, on TriviaQA (500 queries, 66.0\% correct) and NaturalQuestions (500 queries, 19.4\% correct). These values are taken directly from \texttt{h-e1/results/auroc_results.json} and are confirmed by the experiment log (\texttt{h-e1/code/experiment.log}).

\textbf{Table 1: AUROC with 95\% Bootstrap CI — Llama-3-8B-Base}

\begin{table}[htbp]
\centering
\begin{tabular}{lllll}
\toprule
Method & TriviaQA AUROC & TriviaQA 95\% CI & NQ AUROC & NQ 95\% CI \\
\midrule
token_prob & 0.6835 & [0.6361, 0.7332] & 0.6551 & [0.5960, 0.7063] \\
selfcheck_nli & 0.6862 & [0.6362, 0.7340] & 0.4508 & [0.3943, 0.5084] \\
semantic_entropy & 0.4735 & [0.4409, 0.5036] & 0.5524 & [0.5121, 0.5977] \\
kle & 0.2642 & [0.2158, 0.3107] & 0.3753 & [0.3078, 0.4372] \\
selfcheck_bertscore & 0.5000 & [0.5000, 0.5000] & 0.5000 & [0.5000, 0.5000] \\
seps & N/A & N/A & N/A & N/A \\
\bottomrule
\end{tabular}
\end{table}

\textit{SEPs excluded due to zero valid scores on both datasets.}

\begin{figure}[htbp]
\centering
\includegraphics[width=\columnwidth]{figures/fig1_auroc_bar_8b.png}
\caption{AUROC comparison bar chart for all UQ methods on Llama-3-8B-Base}
\label{fig:fig1_auroc_bar_8b}
\end{figure}

\textit{Figure 1: AUROC for each UQ method on TriviaQA and NaturalQuestions (Llama-3-8B-Base). Error bars indicate 95\% bootstrap CI. SEPs omitted due to zero valid scores.}

The existence gate (SE AUROC > token-probability AUROC, 95\% CI of difference excluding zero, on both datasets) fails on both TriviaQA and NaturalQuestions. On TriviaQA, the SE–token-probability difference is −0.210 AUROC points (95\% CI approximately [−0.252, −0.155], excluding zero in the negative direction). On NaturalQuestions, the difference is approximately −0.103 AUROC points.

\begin{figure}[htbp]
\centering
\includegraphics[width=\columnwidth]{figures/fig2_se_tp_difference_8b.png}
\caption{SE minus token-probability AUROC difference with bootstrap CI}
\label{fig:fig2_se_tp_difference_8b}
\end{figure}

\textit{Figure 2: SE minus token-probability AUROC difference with 95\% bootstrap CI, for TriviaQA and NaturalQuestions. Both intervals are entirely negative, excluding zero.}

\subsubsection{5.2 Sampling Degeneracy}

Table 2 reports sampling degeneracy statistics computed by the SE clustering module.

\textbf{Table 2: Sampling Degeneracy Statistics — Llama-3-8B-Base, N=10}

\begin{table}[htbp]
\centering
\begin{tabular}{lllll}
\toprule
Dataset & degenerate_fraction & mean_k & min_k & max_k \\
\midrule
TriviaQA & 0.894 & 9.884 & 7 & 10 \\
NaturalQuestions & 0.848 & 9.796 & 6 & 10 \\
\bottomrule
\end{tabular}
\end{table}

\textit{degenerate_fraction: proportion of queries where all N=10 samples fall into a single semantic cluster (K=1). mean_k: mean number of samples in the dominant cluster per query, out of N=10. Note: mean_k is not the mean number of clusters; for queries with K=1, the dominant cluster contains all 10 samples.}

On TriviaQA, 89.4\% of 500 queries produce K=1: all 10 stochastic samples are assigned to a single semantic equivalence class by the NLI clustering step. For these queries, SE = 0 by definition, regardless of whether the query is answered correctly. The mean_k value of 9.884 indicates that on average across all queries (including the 10.6\% with K>1), the dominant cluster accounts for 98.84\% of the 10 samples. On NaturalQuestions, 84.8\% of queries are degenerate.

The SE mechanism is confirmed to be functioning: the log entry from \texttt{uq_methods.py} at 05:36:52 UTC reads "SE mechanism activated: mean K=9.88 < N=10," confirming that clustering runs and that K < N on every query. The degenerate_fraction captures the degree to which near-trivial clustering is occurring.

\subsubsection{5.3 Method-Specific Findings}

\textbf{Token-probability} achieves AUROC of 0.6835 on TriviaQA and 0.6551 on NaturalQuestions. These values are consistent with the approximately 0.67 AUROC reported for Llama-2-70B in Farquhar et al. (2024). Token-probability is valid under any sampling diversity condition.

\textbf{SelfCheckGPT-NLI} achieves AUROC of 0.6862 on TriviaQA—numerically higher than token-probability—but 0.4508 on NaturalQuestions, below the random-chance baseline. The 95\% CI for NaturalQuestions [0.3943, 0.5084] crosses 0.5, indicating the NQ result is not statistically distinguishable from chance at this sample size. The dataset-dependent behavior of SelfCheckGPT-NLI is not fully explained by the data available; the NaturalQuestions queries have a substantially lower correctness rate (19.4\% vs. 66.0\%), which may affect the method's calibration.

\textbf{KLE} achieves AUROC of 0.2642 on TriviaQA and 0.3753 on NaturalQuestions—substantially below the random-chance baseline of 0.5 on TriviaQA. The 95\% CI [0.2158, 0.3107] for TriviaQA entirely excludes 0.5 in the negative direction. KLE operates on the von Neumann entropy of the pairwise similarity matrix; when 89\% of queries produce rank-1 Laplacian matrices (all samples identical), eigenvalue structure is near-degenerate and the resulting scores are systematically near-zero, potentially with inverted rank ordering relative to true uncertainty.

\textbf{SelfCheckGPT-BERTScore} achieves AUROC of exactly 0.5000 on both datasets, with a zero-width confidence interval [0.5000, 0.5000]. This indicates that the uncertainty score is constant across all queries—a consequence of all samples being identical (BERTScore of identical texts is 1.0), producing a constant score with no rank variation.

\textbf{SEPs} produced zero valid scores on both datasets and are excluded from all analyses. The cause is insufficient probe data at the 500-sample proof-of-concept scale.

\begin{figure}[htbp]
\centering
\includegraphics[width=\columnwidth]{figures/fig3_auroc_all_methods_8b.png}
\caption{All six UQ methods on both datasets}
\label{fig:fig3_auroc_all_methods_8b}
\end{figure}

\textit{Figure 3: AUROC for all evaluated UQ methods on TriviaQA and NaturalQuestions. SelfCheckGPT-BERTScore and SEPs omitted from visualization due to constant scores and null scores respectively.}

\subsubsection{5.4 Gate Evaluation}

The MUST_WORK gate is evaluated against four conditions:

\begin{table}[htbp]
\centering
\begin{tabular}{llll}
\toprule
Condition & Required & Observed & Result \\
\midrule
SE AUROC > token_prob (8B/TriviaQA), CI excl. 0 & SE > TP & SE=0.4735 < TP=0.6835 & FAIL \\
SE AUROC > token_prob (8B/NQ), CI excl. 0 & SE > TP & SE=0.5524 < TP=0.6551 & FAIL \\
SE AUROC > token_prob (70B/TriviaQA), CI excl. 0 & SE > TP & Not available & Unknown \\
SE AUROC > token_prob (70B/NQ), CI excl. 0 & SE > TP & Not available & Unknown \\
\bottomrule
\end{tabular}
\end{table}

The gate fails based on 8B results alone. The 70B experiment did not complete before the reporting cutoff. Because all four conditions must be satisfied and both 8B conditions already fail, 70B results cannot change the gate outcome.

